\topic{先把 64-bit 虚拟地址切成五段}

四级分页使用虚拟地址 bit 47:12 选择四张表，bit 11:0 选择 4 KiB 页内 byte：

\begin{center}
\resizebox{0.98\textwidth}{!}{%
\begin{tikzpicture}[font=\footnotesize\sffamily]
  \foreach \x/\w/\label/\range in {
    0/2.2/符号扩展/63:48,
    2.2/2.2/PML4/47:39,
    4.4/2.2/PDPT/38:30,
    6.6/2.2/PD/29:21,
    8.8/2.2/PT/20:12,
    11.0/2.6/页内偏移/11:0} {
    \draw[BookRule,fill=BookCodeBg] (\x,0) rectangle ++(\w,1.0);
    \node at ({\x+\w/2},0.62) {\label};
    \node at ({\x+\w/2},0.25) {bits \range};
  }
\end{tikzpicture}
}
\end{center}

每个索引有 9 bit，可选 512 项；每项 8 byte，一张表正好是
\[
512\times8=4096\ \text{byte}.
\]
给定表的物理基址 \(B\) 和索引 \(i\)，表项物理地址是：
\[
E=B+8i.
\]
页表遍历不是在虚拟地址上连续“加四次”，而是每一级表项给出下一张表的物理
页帧，再用下一个索引进入那张表。

\topic{canonical 地址先检查高位复制}

当前四级模式使用 bit 47 作为符号位，bit 63:48 必须全部等于 bit 47。

\begin{osgridlongtable}{L{0.29\linewidth}L{0.18\linewidth}L{0.43\linewidth}}
地址 & 是否 canonical & 判断过程 \\ \hline
\endfirsthead
地址 & 是否 canonical & 判断过程 \\ \hline
\endhead
\texttt{0000000000401234} & 是 & bit 47=0，高 16 bit 全 0 \\ \hline
\texttt{00007FFFFFFFFFFF} & 是 & 低 canonical 区最后一个地址 \\ \hline
\texttt{0000800000000000} & 否 & bit 47=1，高 16 bit 却为 0 \\ \hline
\texttt{FFFF800000000000} & 是 & bit 47=1，高 16 bit 全 1 \\ \hline
\texttt{FFFF888012345678} & 是 & 高半区 direct-map 示例 \\ \hline
\end{osgridlongtable}

非 canonical 线性地址不会正常走到某个 PTE。依据触发它的指令类别，处理器
会报告 \#GP 或 \#SS 等异常。映射 API 应在写任何页表之前拒绝该地址，避免
先分配两张表，最后才发现输入根本不属于可翻译集合。

\begin{workedexample}
检查 \texttt{0xFFFF888012345678}。bit 47 位于低 48 bit 的最高位。地址的
高四个十六进制数字是 \texttt{FFFF}，bit 47 也是 1，因此高 16 bit 正是
符号扩展，地址 canonical。

检查 \texttt{0x0000888012345678}。它的 bit 47=1，但高 16 bit 是
\texttt{0000}。正确扩展应以 \texttt{FFFF} 开头，所以该地址非 canonical。
二者低 48 bit 相同，高位检查结果不同。
\end{workedexample}

\topic{算例一：4 KiB 用户代码页}

翻译虚拟地址 \(\texttt{0x0000000000401234}\)。逐级索引和偏移为：

\[
\begin{aligned}
i_4&=(v\gg39)\mathbin{\&}\texttt{1FF}=\texttt{000},\\
i_3&=(v\gg30)\mathbin{\&}\texttt{1FF}=\texttt{000},\\
i_2&=(v\gg21)\mathbin{\&}\texttt{1FF}=\texttt{002},\\
i_1&=(v\gg12)\mathbin{\&}\texttt{1FF}=\texttt{001},\\
offset&=v\mathbin{\&}\texttt{FFF}=\texttt{234}.
\end{aligned}
\]

假设 \(CR3=\texttt{0x00101000}\)，且物理内存中的表项如下：

\begin{osgridlongtable}{L{0.10\linewidth}L{0.27\linewidth}L{0.22\linewidth}L{0.21\linewidth}}
级别 & 表项物理地址 & 表项值 & 解释 \\ \hline
\endfirsthead
级别 & 表项物理地址 & 表项值 & 解释 \\ \hline
\endhead
PML4 & \texttt{101000}+8(0)=\texttt{101000} &
  \texttt{0000000000102007} & 下一表 \texttt{102000}，P/RW/US \\ \hline
PDPT & \texttt{102000}+8(0)=\texttt{102000} &
  \texttt{0000000000103007} & 下一表 \texttt{103000}，P/RW/US \\ \hline
PD & \texttt{103000}+8(2)=\texttt{103010} &
  \texttt{0000000000104007} & 下一表 \texttt{104000}，P/RW/US \\ \hline
PT & \texttt{104000}+8(1)=\texttt{104008} &
  \texttt{0000000012345005} & 页帧 \texttt{12345000}，P/US、只读 \\ \hline
\end{osgridlongtable}

最后物理地址为：
\[
\texttt{0x12345000}+\texttt{0x234}
=\boxed{\texttt{0x12345234}}.
\]
四级都 present；四级都允许 user；叶项 RW=0，所以用户能读和取指，不能写。

\topic{算例二：带 NX 的用户栈页}

翻译 \(\texttt{0x00007FFF12345ABC}\)：
\[
(i_4,i_3,i_2,i_1,offset)
=(\texttt{0FF},\texttt{1FC},\texttt{091},\texttt{145},\texttt{ABC}).
\]

假设前三个表项都是 P/RW/US，叶项为
\(\texttt{0x8000000034567007}\)。bit 63 的 NX=1，页帧基址为
\texttt{0x34567000}，因此：
\[
p=\texttt{0x34567000}+\texttt{0xABC}
=\boxed{\texttt{0x34567ABC}}.
\]

数据读写满足权限；从同一地址取指会触发 protection fault。NX 不改变物理
地址计算，只改变“这次访问能否完成”。

\topic{算例三：高半区内核代码}

翻译 \(\texttt{0xFFFFFFFF80001234}\)：
\[
(i_4,i_3,i_2,i_1,offset)
=(\texttt{1FF},\texttt{1FE},\texttt{000},\texttt{001},\texttt{234}).
\]

假设叶项为 \texttt{0x0000000000200001}：P=1、RW=0、US=0、NX=0。
最终地址：
\[
\texttt{0x00200000}+\texttt{0x234}
=\boxed{\texttt{0x00200234}}.
\]

Ring 0 可读和执行；用户态因 US=0 不能访问。若任一上级项 US=0，即使叶项误设
US=1，用户访问仍失败。权限从每一级共同得到。

\topic{算例四：2 MiB direct-map 大页}

翻译 \(\texttt{0xFFFF8880005ABCDE}\)：
\[
(i_4,i_3,i_2)=(\texttt{111},\texttt{000},\texttt{002}),\qquad
offset_{2M}=\texttt{1ABCDE}.
\]

PDE 值为 \texttt{0x8000000000400083}。P=1、RW=1、PS=1、NX=1，物理大页
基址是 \texttt{0x00400000}。PS=1 使遍历在 PD 提前结束，不再读取 PT：
\[
\texttt{0x00400000}+\texttt{0x1ABCDE}
=\boxed{\texttt{0x005ABCDE}}.
\]

\begin{center}
\begin{tikzpicture}[font=\footnotesize\sffamily,node distance=9mm and 13mm]
  \node[os state] (va) {虚拟地址};
  \node[os block,right=of va] (pml4) {PML4E};
  \node[os block,right=of pml4] (pdpt) {PDPTE};
  \node[os block,right=of pdpt] (pde) {PDE\\PS=1};
  \node[os state,below=of pde] (base) {2 MiB 页帧基址};
  \node[os state,left=of base] (offset) {低 21 bit 偏移};
  \node[os state,left=of offset] (pa) {物理地址};
  \draw[os arrow] (va) -- (pml4);
  \draw[os arrow] (pml4) -- (pdpt);
  \draw[os arrow] (pdpt) -- (pde);
  \draw[os arrow] (pde) -- (base);
  \draw[os arrow] (base) -- (offset);
  \draw[os arrow] (offset) -- (pa);
\end{tikzpicture}
\end{center}
\diagramnote{4 KiB 路径使用 12-bit 偏移并访问 PT；2 MiB 路径使用 21-bit
偏移并在 PDE.PS=1 处结束。}

\topic{权限逐级合成，而不是只读叶项}

\begin{osgridlongtable}{L{0.20\linewidth}L{0.24\linewidth}L{0.26\linewidth}L{0.20\linewidth}}
访问条件 & 上级项 & 叶项 & 有效结果 \\ \hline
\endfirsthead
访问条件 & 上级项 & 叶项 & 有效结果 \\ \hline
\endhead
present & 每级 P=1 & P=1 & 才能完成翻译 \\ \hline
用户访问 & 每级 US=1 & US=1 & Ring 3 才可访问 \\ \hline
写访问 & 每级 RW=1 & RW=1 & 页面才可写 \\ \hline
取指 & 每级 NX=0 & NX=0 & 页面才可执行 \\ \hline
\end{osgridlongtable}

有效 RW 和 US 可理解为沿路径做逻辑 AND；有效 NX 可理解为沿路径做逻辑 OR。
任何一级 NX=1 都禁止从该页取指。

\begin{workedexample}
四级 RW 依次为 \(1,1,0,1\)，US 为 \(1,1,1,1\)，NX 为
\(0,0,0,0\)。用户读成功，用户写失败，用户取指成功。PD 的 RW=0 已经限制
整个 2 MiB 范围，叶项 RW=1 不能重新放宽。

若把叶项 NX 改为 1，地址翻译仍得到同一物理地址，但取指失败。权限错误与
not-present 错误必须从页故障错误码区分。
\end{workedexample}

\topic{CR0.WP 决定内核写入是否服从只读页}

当 \(CR0.WP=1\) 时，supervisor 写入也要满足有效 RW。内核
\texttt{.text} 和 \texttt{.rodata} 才能依靠页表阻止意外写入。若 WP=0，
Ring 0 可能写入 supervisor 只读页，代码中的越界写会悄悄修改指令或常量。

启用顺序必须是：

\begin{enumerate}[leftmargin=2.2em]
  \item 先为当前 RIP、RSP、IDT 和数据建立正确映射；
  \item 设置 EFER.NXE，使 NX 位获得定义的执行禁止语义；
  \item 设置 CR0.WP；
  \item 切换到候选 CR3；
  \item 用一次故意写只读页确认 \#PF 的 CR2 和错误码。
\end{enumerate}

先切 CR3 再补当前栈，会在下一次压栈时故障；异常入口若也未映射，处理器无法
保存故障现场，可能继续演变为 double fault 和 triple fault。

\topic{页故障错误码要和 CR2 一起读}

\begin{osgridlongtable}{L{0.12\linewidth}L{0.17\linewidth}L{0.61\linewidth}}
bit & 名称 & 为 1 时的含义 \\ \hline
\endfirsthead
bit & 名称 & 为 1 时的含义 \\ \hline
\endhead
0 & P & 保护检查失败；为 0 时通常是页不存在 \\ \hline
1 & W/R & 访问是写；为 0 时是读 \\ \hline
2 & U/S & 访问来自用户态；为 0 时来自 supervisor \\ \hline
3 & RSVD & 页表项的保留位组合不合法 \\ \hline
4 & I/D & 访问是取指，需启用相应架构能力 \\ \hline
5 & PK & 保护键检查失败 \\ \hline
6 & SS & shadow-stack 访问失败 \\ \hline
\end{osgridlongtable}

\begin{osgridlongtable}{L{0.25\linewidth}L{0.18\linewidth}L{0.47\linewidth}}
场景 & 错误码 & 解码 \\ \hline
\endfirsthead
场景 & 错误码 & 解码 \\ \hline
\endhead
用户读取 not-present 页 & \texttt{0x4} & P=0，read，user \\ \hline
内核写 present 只读页 & \texttt{0x3} & P=1，write，supervisor \\ \hline
用户写 present 只读页 & \texttt{0x7} & P=1，write，user \\ \hline
用户从 NX 页取指 & \texttt{0x15} & P=1，user，instruction fetch \\ \hline
页表保留位错误 & \texttt{0x9} 或更多位 & P=1，RSVD=1，其余位随访问而定 \\ \hline
\end{osgridlongtable}

CR2 保存引发最近一次页故障的线性地址，错误码说明访问类型。异常处理程序应在
允许再次触发 \#PF 的操作前保存 CR2，否则嵌套故障会覆盖它。CR3 给出当前
地址空间的根表物理地址；启用 PCID 时低位还携带标识，取得根表基址时要按
架构掩码清除非地址位。

\begin{workedexample}
Ring 3 执行
\code{mov byte [0x0000000000402000], 1}。目标页 present、US=1，但 PD 的
RW=0，CR0.WP=1。处理器先算出线性地址，再完成四级遍历；最终物理页帧即使
已经算出，写访问仍不能提交。

\begin{osgridlongtable}{L{0.12\linewidth}L{0.30\linewidth}L{0.48\linewidth}}
时刻 & CPU 内部动作 & 可观察状态 \\ \hline
\endfirsthead
时刻 & CPU 内部动作 & 可观察状态 \\ \hline
\endhead
t0 & 形成线性地址 \texttt{402000} & RIP 指向 faulting \code{mov} \\ \hline
t1 & TLB miss，按 CR3 逐级读表 & PML4E、PDPTE、PDE、PTE 都 P=1 \\ \hline
t2 & 合成权限 & US=1，RW=0，所以用户写失败 \\ \hline
t3 & 进入 \#PF & CR2=\texttt{402000}，error=\texttt{07} \\ \hline
t4 & 异常入口保存现场 & 保存的 RIP 仍指向原 \code{mov} \\ \hline
t5 & 处理程序修复 PTE 并 \code{invlpg} & 旧 TLB 项不再可用 \\ \hline
t6 & \code{iretq} 后重试原指令 & 第二次遍历允许写，byte 才变为 1 \\ \hline
\end{osgridlongtable}

错误码 \texttt{07} 的二进制低三位是 \texttt{111}：present 页面上的写访问，
访问来自 user。CR2 回答“哪个线性地址”，错误码回答“怎样访问时失败”，保存
的 RIP 回答“哪条指令要重试”。
\end{workedexample}

\topic{先把 CR3 中的根地址和上下文编号分开}

假设调试器显示 \(\mathrm{CR3}=\texttt{0x0000000000101007}\)，并且
CR4.PCIDE=1。低 12 bit 是 PCID，页表根基址来自其余地址位：
\[
\begin{aligned}
root&=\texttt{0x101007}\mathbin{\&}
      \texttt{0x000FFFFFFFFFF000}
      =\texttt{0x101000},\\
pcid&=\texttt{0x101007}\mathbin{\&}\texttt{0xFFF}
      =\texttt{0x007}.
\end{aligned}
\]

\begin{osgridlongtable}{L{0.22\linewidth}L{0.27\linewidth}L{0.41\linewidth}}
条件 & CR3 低位怎样解释 & 调试时先做什么 \\ \hline
\endfirsthead
条件 & CR3 低位怎样解释 & 调试时先做什么 \\ \hline
\endhead
CR4.PCIDE=0 & PWT、PCD 等控制位，不是 PCID & 查当前分页模式，再取根地址 \\ \hline
CR4.PCIDE=1 & bits 11:0 是 PCID & 同时记录 root 与 PCID \\ \hline
切换 CR3 & 选择新的根表和翻译上下文 & 不把 CR3 数值当普通虚拟指针 \\ \hline
\end{osgridlongtable}

书中手算用固定物理地址掩码帮助观察；正式实现还要根据 CPUID 报告的物理地址
宽度处理地址位和保留位。直接把 \texttt{\symbol{36}cr3} 强制转换成 C++ 指针会让 CPU
用当前虚拟地址空间解释一个物理地址，二者通常不是同一件事。

\topic{TLB 让“内存里的新 PTE”不等于“CPU 已使用新翻译”}

第一次访问后，CPU 可把虚拟页、物理页和权限缓存在 TLB。软件把 PTE 从可写
改为只读后，旧 TLB 项可能仍允许写。单页修改后可执行
\code{invlpg [virtual_address]}；切换地址空间常通过写 CR3 改变当前翻译
上下文。

\begin{pdfdiagram}
  \node[os state] (access1) {第一次访问};
  \node[os block,right=of access1] (walk) {四级 page walk};
  \node[os state,right=of walk] (tlb) {TLB 保存翻译和权限};
  \node[os state,below=of walk] (pte) {软件修改 PTE};
  \node[os block,left=of pte] (stale) {未刷新\\继续命中旧项};
  \node[os block,right=of pte] (flush) {\code{invlpg}\\下一次重新遍历};
  \draw[os arrow] (access1) -- (walk);
  \draw[os arrow] (walk) -- (tlb);
  \draw[os arrow] (tlb) -- (stale);
  \draw[os arrow] (pte) -- (stale);
  \draw[os arrow] (pte) -- (flush);
\end{pdfdiagram}

多核机器的 TLB 属于各个核心。取消映射后，内核要通知曾运行该地址空间的其他
核心，让它们失效并确认，然后才能把物理页交给别的对象。否则旧核心仍可能
通过陈旧翻译写入已经复用的页。

\begin{experimentbox}{用同一个 store 比较刷新前后}
在测试专用页 \code{test_page} 上按下面顺序操作：

\begin{enumerate}[leftmargin=2.2em]
  \item 先写一次 \code{test_page[0]=0x11}，让当前 CPU 建立 TLB 项；
  \item 在 GDB 中暂停，找到叶 PTE，清除 RW，但保持 P=1；
  \item 不执行失效，恢复并再次写 \code{0x22}，记录它是否完成；
  \item 再次暂停，对同一线性地址执行 \code{invlpg}；
  \item 第三次写 \code{0x33}，应进入 \#PF，并记录 CR2 和错误码；
  \item 恢复 PTE.RW，执行 \code{invlpg}，确认重试后写入成功。
\end{enumerate}

第 3 步不能作为跨处理器型号的稳定断言：没有按架构要求失效时，软件已经
失去可靠语义，CPU 可能继续使用缓存项。真正的检查点是第 4--5 步：刷新以后，
下一次访问必须重新看到只读权限。
\end{experimentbox}

\topic{运行页表遍历器，而不只在纸上看伪代码}

\filepath{examples/page\_table\_walker/} 是一个可以单独编译的 C++20
程序。它把一段数组当作物理内存中的页表镜像，从 CR3 指向的根表开始读取。
每读一级，它都保存表项物理地址、原始值和累计后的 RW、US、NX。遍历结束
以后，再根据访问类型、特权级和 CR0.WP 判断这次访问能否完成。

\lstinputlisting[
  language=C++,
  firstline=178,
  lastline=221,
  caption={先检查 canonical 地址，再逐级读取并累计权限}
]{../../examples/page_table_walker/source/page_table_walker.cpp}

4 KiB 路径走完 PT，2 MiB 路径在 PDE.PS=1 时提前结束。地址算完以后仍不能
立刻返回成功，因为同一个物理地址可能允许读、禁止写，或允许数据访问、禁止
取指：

\lstinputlisting[
  language=C++,
  firstline=223,
  lastline=255,
  caption={大页偏移、最终权限检查和页故障错误码}
]{../../examples/page_table_walker/source/page_table_walker.cpp}

\begin{experimentbox}{编译并运行五条遍历路径}
在书稿目录运行：

\begin{lstlisting}[language=bash]
python3 examples/page_table_walker/verify.py
\end{lstlisting}

脚本以 \code{-std=c++20 -Wall -Wextra -Wpedantic -Werror} 编译声明、实现和
驱动程序，然后运行五个场景：

\begin{osgridlongtable}{L{0.25\linewidth}L{0.23\linewidth}L{0.18\linewidth}L{0.24\linewidth}}
场景 & 结果 & 错误码 & 要检查的中间状态 \\ \hline
\endfirsthead
场景 & 结果 & 错误码 & 要检查的中间状态 \\ \hline
\endhead
4 KiB 用户读 & \code{translated} & 0 & 四级 entry-address 与最终偏移 \\ \hline
2 MiB supervisor 读 & \code{translated} & 0 & 只访问三级，offset 为 21 bit \\ \hline
用户从 NX 页取指 & \code{execute-denied} & \texttt{15} & 地址算出后权限失败 \\ \hline
父级只读、用户写 & \code{write-denied} & \texttt{07} & 叶 RW=1 也不能放宽父级 \\ \hline
用户读 not-present 页 & \code{not-present} & \texttt{04} & 停在第一个 P=0 表项 \\ \hline
\end{osgridlongtable}

不要只看脚本最后一行“通过”。把 \code{scenario=4k-read} 下的四个
\code{entry-address} 与算例一逐项对照；随后把驱动中的叶项 US 清零，重新
运行，确认物理地址没有变化，但状态改为 \code{user-denied}。
\end{experimentbox}

这个程序仍然刻意省略 1 GiB 页、保留位检查和真实物理地址宽度。它只查询
已有页表，不为缺失层分配页帧。正式内核还必须区分“查询映射”和“创建映射”，
创建路径失败时还要撤销已经写入的表项。

\topic{第四个页帧申请失败时怎样回到原状态}

假设要映射四个虚拟页，前三个已经取得物理页并建立 PTE，第四次物理页申请
失败。具体恢复动作是：

\begin{enumerate}[leftmargin=2.2em]
  \item 把已写入的三个叶 PTE 按反向顺序清零；
  \item 若本次新建的 PT 已无有效叶项，从 PDE 断开它并释放 PT 页帧；
  \item 对本次新建且已经变空的 PD、PDPT 采用同样处理；
  \item 释放前三个数据页帧；
  \item 若修改的是当前活动地址空间，对三个虚拟页执行失效；候选 CR3 尚未
    激活时没有对应的活动 TLB 项；
  \item 返回申请失败，调用者看到的原映射保持不变。
\end{enumerate}

\begin{osgridlongtable}{L{0.16\linewidth}L{0.22\linewidth}L{0.24\linewidth}L{0.28\linewidth}}
时刻 & 已取数据页 & 已写 PTE & 下一动作 \\ \hline
\endfirsthead
时刻 & 已取数据页 & 已写 PTE & 下一动作 \\ \hline
\endhead
开始 & 0 & 0 & 记录原表项 \\ \hline
页 1 成功 & 1 & 1 & 继续页 2 \\ \hline
页 2 成功 & 2 & 2 & 继续页 3 \\ \hline
页 3 成功 & 3 & 3 & 请求页 4 \\ \hline
页 4 失败 & 3 & 3 & 清 PTE 3、2、1 \\ \hline
恢复完成 & 0 & 0 & 释放变空的新表并返回失败 \\ \hline
\end{osgridlongtable}

如果只释放数据页却保留三个 present PTE，后续分配器会把这些页给新对象，旧
虚拟地址仍能访问它们。若只清 PTE 不释放新建中间表，重复失败会逐步泄漏页表
页。所谓失败原子性在这里就是上述每一个可检查动作。

\topic{用 QEMU 和 GDB 逐级查看页表}

\begin{experimentbox}{从 CR3 找到一个 PTE}
以带 GDB stub 的 QEMU 启动系统，在内核进入稳定日志点后暂停。

\begin{enumerate}[leftmargin=2.2em]
  \item 在 GDB 中执行
    \textcolor{BookBlue}{\texttt{p/x \symbol{36}cr3}}，记录根表物理地址；
  \item 选取日志中已知的虚拟地址，手算四级索引；
  \item 使用 QEMU monitor 的物理内存查看命令
    \textcolor{BookBlue}{\texttt{xp/4gx physical\_address}} 读取根表项；
  \item 清除表项低 12 位得到下一级表物理基址，重复到叶项；
  \item 把页帧基址与偏移相加，再用 monitor 查看该物理地址的 byte；
  \item 用 \textcolor{BookBlue}{\texttt{monitor info tlb}} 或
    \textcolor{BookBlue}{\texttt{monitor info mem}} 辅助观察，
    但最终索引与地址仍要在纸上独立计算。
\end{enumerate}

命令名称可能随 QEMU monitor 前端略有差异；关键记录是 CR3、四个索引、每级
表项原值、最终页帧和偏移。

若使用算例一的页表镜像，完整会话应能整理成下面这张表。命令中的地址是物理
地址，不是让 GDB 按当前虚拟地址空间解引用：

\begin{osgridlongtable}{L{0.16\linewidth}L{0.37\linewidth}L{0.37\linewidth}}
动作 & 命令或手算 & 预期值 \\ \hline
\endfirsthead
动作 & 命令或手算 & 预期值 \\ \hline
\endhead
读根表 & \textcolor{BookBlue}{\texttt{p/x \symbol{36}cr3}} &
  \texttt{0x101000} \\ \hline
读 PML4E & \textcolor{BookBlue}{\texttt{monitor xp /1gx 0x101000}} &
  \texttt{0x102007} \\ \hline
读 PDPTE & \textcolor{BookBlue}{\texttt{monitor xp /1gx 0x102000}} &
  \texttt{0x103007} \\ \hline
读 PDE & \textcolor{BookBlue}{\texttt{monitor xp /1gx 0x103010}} &
  \texttt{0x104007} \\ \hline
读 PTE & \textcolor{BookBlue}{\texttt{monitor xp /1gx 0x104008}} &
  \texttt{0x12345005} \\ \hline
读目标 byte & \textcolor{BookBlue}{\texttt{monitor xp /1bx 0x12345234}} &
  与虚拟地址读取相同 \\ \hline
\end{osgridlongtable}

若第二级就读到 0，不要继续用书中的下一级地址。先确认当前 CR3、所选虚拟
地址和索引；每一级的新基址都必须来自刚刚读到的表项。若 monitor 不提供
\code{info tlb}，逐级物理内存读取仍然足以验证页表；TLB 是否刷新则用前面的
权限故障实验验证。
\end{experimentbox}

\begin{faultinjectionbox}{清除一个父级 RW 并写入子页}
选择一张测试专用页表，把目标路径中 PD 的 RW 清零，叶 PTE 保持 RW=1。
执行 \code{invlpg} 后从 Ring 0 写该页，并确保 CR0.WP=1。

预期得到 \#PF，CR2 等于测试虚拟地址，错误码为 \texttt{0x3}。若写入成功，
先检查 WP；若错误码为 \texttt{0x2}，说明某级 P=0，测试没有到达只读权限
检查。恢复 PD 原值并再次失效后，写入应成功。
\end{faultinjectionbox}

\begin{faultinjectionbox}{切换到缺少当前栈的候选 CR3}
只在专用故障镜像中进行。候选表映射当前 RIP 和异常入口，但故意不映射当前
RSP 所在页。写 CR3 后执行一次 \code{push}。预期 CR2 落在栈顶下方，错误码
P=0、W/R=1。若异常入口所需栈也缺失，可能看不到普通 \#PF 日志并发生重启，
这正说明激活前必须同时检查代码、栈和异常路径。
\end{faultinjectionbox}

\begin{pitfallbox}
\begin{itemize}[leftmargin=2.2em]
  \item CR3 给出根表物理地址，不是一个可直接解引用的普通 C++ 指针；
  \item 2 MiB 页在 PDE.PS=1 处结束，偏移是 21 bit；
  \item 叶项 RW=1 不能放宽父级 RW=0；
  \item 修改 PTE 后要处理 TLB；只看内存中的新值可能误判；
  \item CR2 是故障地址，错误码是访问类别，两者缺一不可；
  \item 分配失败要清已写 PTE、释放数据页和本次新建的空表。
\end{itemize}
\end{pitfallbox}

\topic{练习：手算、遍历、改错}

\begin{enumerate}[leftmargin=2.2em]
  \item 求 \texttt{0x00007FFF12345ABC} 的四级索引和 4 KiB 偏移。
  \item 2 MiB 页基址为 \texttt{0x02000000}，虚拟地址页内偏移为
    \texttt{0x1F0123}。求物理地址。
  \item 四级 US 为 \(1,1,0,1\)，RW 为 \(1,1,1,1\)，NX 为
    \(0,0,0,0\)。用户读、内核写、用户取指分别能否完成？
  \item 解码错误码 \texttt{0x15}。CR2 在异常处理程序中提供什么信息？
  \item 映射 8 页时第 6 个数据页申请失败。列出要清理的叶项数量和数据页
    数量，并说明何时需要释放中间表。
\end{enumerate}

\begin{answerbox}
\begin{enumerate}[leftmargin=2.2em]
  \item \((\texttt{0FF},\texttt{1FC},\texttt{091},
    \texttt{145})\)，偏移 \texttt{ABC}；
  \item \texttt{0x02000000}+\texttt{0x1F0123}
    \(=\texttt{0x021F0123}\)；
  \item 用户读失败，因为父级 US=0；内核写可完成，所有 RW=1；用户取指也因
    US=0 失败；
  \item \texttt{0x15} 表示 present 页面上的用户取指违反保护，常见原因是
    NX。CR2 给出发生这次取指的线性地址；
  \item 已成功 5 页，所以清 5 个叶项并释放 5 个数据页。只释放本次创建且
    清理后没有任何有效子项的中间表，不能释放原先已有或仍含其他映射的表。
\end{enumerate}
\end{answerbox}
